% Generated draft: requires mathematical and source review.
\documentclass[11pt]{article}
\usepackage[margin=25mm]{geometry}
\usepackage{fontspec}
\setmainfont{Latin Modern Roman}
\usepackage{amsmath,amssymb,mathtools,unicode-math}
\setmathfont{Latin Modern Math}
\usepackage{enumitem,needspace,xurl,hyperref}
\setlength{\parindent}{0pt}
\setlength{\parskip}{5pt}
\setlength{\emergencystretch}{3em}
\newcommand{\N}{\mathbb{N}}
\newcommand{\Z}{\mathbb{Z}}
\newcommand{\Q}{\mathbb{Q}}
\newcommand{\R}{\mathbb{R}}
\newcommand{\C}{\mathbb{C}}
\newcommand{\F}{\mathbb{F}}
\newcommand{\A}{\mathbb{A}}
\newcommand{\PP}{\mathbb{P}}
\newcommand{\E}{\mathbb{E}}
\newcommand{\eps}{\varepsilon}
\newcommand{\dd}{\,\mathrm{d}}
\begin{document}
\section*{UnsolvedMath ID: 30000077}
\textbf{Source problem code:} catalogue entry
\subsection{Definitions and mathematical statement}
Let $S=k[x_1,\ldots,x_m]$ be a standard graded polynomial ring over a field $k$, with its given variables, and let $I\subsetneq S$ be a homogeneous ideal in the source's intended class of ``classical'' ideals. For a monomial term order $\prec$, let
\[
 H=\operatorname{in}_{\prec}(I)
\]
be the ideal generated by leading monomials in those same coordinates. For a finitely generated graded $S$-module $M$, put
\[
 \beta^S_{i,j}(M)=\dim_k\operatorname{Tor}^S_i(M,k)_j.
\]
Here the bidegree $(i,j)$ consists of homological degree and internal polynomial degree; it does not presume a separate bigrading on $S$.

The question is whether every such classical ideal admits a term order in the fixed coordinates for which
\[
 \beta^S_{i,j}(S/I)=\beta^S_{i,j}(S/H)
 \qquad\text{for all }i\ge0,\ j\in\Z.
\]
Equivalently one may compare the ideals' graded Betti numbers with the usual homological shift, provided the same convention is used for both. Coordinate changes are not part of the allowed existence assertion. Equality of Hilbert functions alone, or an initial ideal generated by quadrics alone, is weaker than this requirement.

The extract does not define the full family meant by ``classical'', specify the field or characteristic, or state whether any generalized weight-initial constructions beyond monomial term orders are allowed. These are source-scope qualifications; the question is not expanded to all homogeneous ideals.

The literature assessment checked on 21 August 2026 says the universal Betti-preserving initial-ideal question remains unsettled. It notes determinantal examples where quadratic initial ideals acquire nonlinear syzygies, and positive constructions for several classical families in later work. A bad initial ideal for one order does not disprove the existence of a different Betti-preserving order. This is dated supplied triage requiring independent verification.
\subsection{Short English statement}
Can each intended classical ideal acquire a monomial initial ideal with the same entire graded Betti table while keeping its coordinates fixed? Known successes in some families and failures for particular quadratic initial ideals do not settle that existence question.
\subsection{Sources}
\begin{itemize}
\item \emph{Combinatorial Commutative Algebra} (2005), Oberwolfach Reports, DOI 10.4171/owr/2004/32, pp. 1703--1746, \url{https://doi.org/10.4171/owr/2004/32}; proposed year supplied as 2004.
\item Literature assessment checked 21 August 2026: Conca--Hosten--Thomas, \emph{Nice Initial Complexes}, \url{https://arxiv.org/abs/math/0512283}; ``Original Oberwolfach report,'' \url{https://doi.org/10.4171/owr/2004/32}.
\item Local UnsolvedMath registry, permanent ID $30000077$, exact source code \texttt{catalogue entry}, created 25 August 2026 and updated 16 September 2026. The clean statement is labelled corrected and verified against the report.
\end{itemize}
\end{document}
